%% ICSOFT 2026 -- Abstracts Track
%% Template: SCITEPRESS (SCITEPRESS.sty, article.cls, apalike.sty/bst in this dir)
%%
%% Status: DRAFT for co-author review (sent to Philippsen 2026-05-05)
%%   Pending: (1) English translation
%%            (2) Veronika Schwarz confirmation
%%            (3) Trim to ~350 words for 1-page fit
%%            (4) Verify numbers: 92 modules software-ekg-7, 4.038/3.585 Minecraft

\documentclass[a4paper,twoside]{article}

\usepackage{epsfig}
\usepackage{subcaption}
\usepackage{amssymb}
\usepackage{amsmath}
\usepackage{amsthm}
\usepackage{multicol}
\usepackage{pslatex}
\usepackage{apalike}
\usepackage{caption}
\usepackage[ruled]{algorithm2e}
\usepackage{SCITEPRESS}
\usepackage[utf8]{inputenc}
\usepackage[T1]{fontenc}

%% ---------------------------------------------------------------
\title{Korrekte Layer-Berechnung f\"{u}r hierarchische
       Architekturvisualisierungen
       \subtitle{Constraints, Invarianten und Pitfalls
       \\ \normalsize\textit{[Entwurf zur Abstimmung --- bitte um R\"{u}ckmeldung bis 17.~Mai~2026]}}}

\author{
  \authorname{Johannes Weigend}
  \affiliation{Technische Hochschule Rosenheim / Weigend AM GmbH \& Co.KG, Germany}
  \email{johannes.weigend@gmail.com}
  \authorname{Veronika Schwarz}
  \affiliation{[Affiliation bitte erg\"{a}nzen]}%  % formerly published as Veronika Dashuber
  \email{[E-Mail bitte erg\"{a}nzen]}
  \authorname{Michael Philippsen}
  \affiliation{Friedrich-Alexander-Universit\"{a}t Erlangen-N\"{u}rnberg, Germany}
  \email{[E-Mail bitte erg\"{a}nzen]}
}

\keywords{software architecture visualization, layout invariants,
  layered software city, strongly connected components,
  dependency analysis, tool demonstration}

\abstract{Geschichtete Architekturvisualisierungen dr\"{u}cken Abh\"{a}ngigkeitsstruktur
durch vertikale Anordnung aus: wer von anderen abh\"{a}ngt, steht h\"{o}her.
Ein Package auf einer niedrigeren Schicht als seine eigenen Klassen,
zwei zyklisch abh\"{a}ngige Pakete auf unterschiedlichen Ebenen, eine
regul\"{a}re Abh\"{a}ngigkeit die entgegen der Schichtrichtung zeigt~---
solche Fehler in der Layer-Berechnung sehen plausibel aus, t\"{a}uschen
aber den Betrachter. Das Prinzip klingt selbstverst\"{a}ndlich: die
Korrektheit eines Architekturlayouts sollte das Werkzeug garantieren,
nicht der Nutzer. Wie schwer das tats\"{a}chlich ist, zeigt erst die
Implementierung.
\par\smallskip
Unser 2021 bei IVAPP ausgezeichnetes Paper \cite{Dashuber2021}
f\"{u}hrte ein Schichtlayout ein, dessen Grundprinzip einfach ist: jede
Klasse erh\"{a}lt ein h\"{o}heres Level als die Klassen von denen sie
abh\"{a}ngt. Die X/Y-Position von Geb\"{a}uden auf dem Stadtgrundriss
folgt dieser Ordnung direkt~--- foundationale Klassen stehen vorne/unten,
hochrangige Module hinten/oben. Geb\"{a}udeh\"{o}hen kodieren orthogonal
dazu Metriken wie Komplexit\"{a}t oder Zeilenzahl.
Dieser Beitrag stellt eine quelloffene Referenzimplementierung (Java,
Apache~2.0) vor sowie einen \textbf{Layout-Invarianten-Checker} der
vier formal spezifizierte Nachbedingungen nach jedem Analyselauf
verifiziert.
\par\smallskip
Die Pipeline selbst ist nicht trivial. Ein korrektes Schichtlayout
erfordert die L\"{o}sung dreier interdependenter Constraint-Systeme~---
Klassen-Ordnung, Package-Hierarchie und SCC-Gleichheit. Ein naiver
Single-Pass-Algorithmus liefert bei zyklischen Graphen
nicht-deterministische Ergebnisse. Der konzeptionelle Kern der
L\"{o}sung: die Pipeline berechnet \textbf{zwei semantisch verschiedene
Level-Konzepte}. Phase~1 berechnet die \emph{architektonische Tiefe}
als l\"{a}ngsten Pfad im globalen, zyklenfreien SCC-DAG. Phase~2
berechnet die \emph{lokale Layout-Position} innerhalb eines Packages,
beginnend bei~0, ausschlie{\ss}lich auf Geschwister-Abh\"{a}ngigkeiten.
Stark zyklischer Code~--- Minecraft Forge: eine SCC aus 4.038~Klassen
mit 3.585~R\"{u}ckkanten~--- wird durch einen heuristischen SCC-Breaker
\cite{Eades1993} zu einer nutzbaren Hierarchie aufgel\"{o}st.
\par\smallskip
Der Invarianten-Checker (R1--R5) unterscheidet pr\"{a}zise zwischen
Algorithmenfehlern und tolerierten Heuristik-Verletzungen. Die
Entwicklung dieser Regeln war selbst nicht fehlerfrei: eine fr\"{u}he
Variante von R2 erwies sich als zu grob und musste verfeinert werden.
Der Checker validiert damit nicht nur die Pipeline, sondern auch die
Pr\"{a}zision der Regeln selbst~--- ein doppelter Boden den kein
manueller Review ersetzen kann. Dieselben vier Regeln wurden portiert
auf eine unabh\"{a}ngige 3D-Software-City in Unity/.NET~--- mit
byte-identischen Befunden auf gemeinsamen Eingaben.
Ein R2-Befund in der 92-Modul-Codebasis \emph{software-ekg-7}
identifizierte einen fehlenden Package-SCC-Equalisierungsschritt und
l\"{o}ste direkt einen Pipeline-Fix aus. Nach unserem Kenntnisstand hat
keine fr\"{u}here geschichtete Architekturvisualisierung diesen Grad an
Selbstvalidierung realisiert.}
%% ---------------------------------------------------------------

\begin{document}

\maketitle

\vfill

\bibliographystyle{apalike}
{\small
\bibliography{icsoft2026}}

%% ===================================================================
\appendix
\onecolumn

\section*{ANHANG A --- Warum Single-Pass nicht funktioniert:
          eine Analyse der Fehlermodi}

\subsection*{Das Grundproblem: drei interagierende Constraint-Systeme}

Ein korrektes Schichtlayout muss gleichzeitig drei Constraints erf\"{u}llen:

\begin{enumerate}
  \item \textbf{Klassen-Constraint}: Wenn $A \rightarrow B$ ($A$ h\"{a}ngt von $B$ ab),
        dann muss $\text{level}(A) > \text{level}(B)$.
  \item \textbf{Hierarchie-Constraint (R3)}: Ein Package muss auf mindestens dem Level
        seines h\"{o}chsten Inhalts liegen.
  \item \textbf{SCC-Constraint (R2)}: Alle Klassen in einer zyklischen Abh\"{a}ngigkeit
        bekommen dasselbe Level; alle Packages die eine Package-SCC bilden ebenso.
\end{enumerate}

Diese drei Constraints sind einzeln l\"{o}sbar. Zusammen erzeugen sie ein
\textbf{zirkul\"{a}res Abh\"{a}ngigkeitssystem zwischen den Ebenen der Berechnung selbst}
--- nicht zwischen den Klassen.

\subsection*{Die konkreten Fehlermodi eines Single-Pass-Ansatzes}

\textbf{Fehlermodus 1: Retroaktive Korrektur.}
Der Algorithmus iteriert \"{u}ber die Klassen und vergibt Level.
Er sieht Klasse $A$ (in Package $P_1$), berechnet $\text{level}(A)=3$,
setzt damit $\text{level}(P_1)=3$. Weiter trifft er Klasse $B$ (in $P_2$) und
stellt fest: $A$ und $B$ sind in einer SCC, m\"{u}ssen gleiches Level bekommen,
jetzt $\text{level}(A)=\text{level}(B)=5$. Damit muss $\text{level}(P_1)$
neu berechnet werden --- aber $P_1$ und $P_3$ wurden bereits abgeschlossen.
Die SCC-Mitgliedschaft ist erst nach vollst\"{a}ndiger Graphtraversierung bekannt.

\textbf{Fehlermodus 2: Verwobene Zykluserkennung erzeugt classpath-abh\"{a}ngige
Ergebnisse.}
Klassen kommen aus JARs. Die Reihenfolge folgt der JAR-Packreihenfolge, die
abh\"{a}ngt von Build-Tool (Maven vs.\ Gradle), Modul-Reihenfolge und
Dateisystem-Sortierung --- keine davon codiert Abh\"{a}ngigkeitsrichtung.
Derselbe Quellcode liefert so unterschiedliche Level-Layouts je nach
Build-Konfiguration. Die Fehler sind \textbf{subtil}: Das Layout sieht plausibel aus,
Pfeile zeigen grob in die richtige Richtung, die pr\"{a}zise Hierarchieinformation
ist korrumpiert aber nicht sofort sichtbar.

\textbf{Fehlermodus 3: Gegenseitiger Verweis zwischen Package- und Klassen-Level.}
$\text{level}(P_1) = \max(\text{level}(c) \mid c \in \text{subtree}(P_1))$.
Welches Level hat die h\"{o}chste Klasse? Das h\"{a}ngt von anderen Klassen ab,
die in Packages liegen, deren Level davon abh\"{a}ngt\ldots\ Dieser zirkul\"{a}re
Verweis macht einen gemeinsamen Pass \"{u}ber Klassen und Packages inh\"{a}rent instabil.

\textbf{Fehlermodus 4: Der gro{\ss}e See.}
Ohne Mechanismus zum Aufbrechen gro{\ss}er SCCs landen alle Mitglieder
zwangsl\"{a}ufig auf demselben Level. Bei Minecraft Forge --- eine SCC aus
4.038~Klassen --- liegen \"{u}ber 40\,\% aller Klassen auf Level~0: ein riesiger,
flacher See am unteren Rand des Layouts. Der heuristische SCC-Breaker
(In-/Out-Degree-Rankscore, verwandt mit \cite{Eades1993}) l\"{o}st dieses Problem.
Entfernte Kanten werden als Verletzungen (rot) visualisiert, nicht ignoriert.

\subsection*{Formaler Beweis: Warum ein naiver rekursiver Algorithmus versagt}

Der intuitivste Ansatz ist ein rekursiver DFS mit Memoization:

\begin{verbatim}
level(A):
  if A already computed: return level(A)
  if A has no dependencies: return 0
  return max(level(dep) for dep in A.dependencies) + 1
\end{verbatim}

F\"{u}r DAGs korrekt. F\"{u}r zyklische Graphen versagt er strukturell.

\textbf{Gegenbeispiel} (zwei verschachtelte SCCs):
$\text{SCC}_1 = \{A,B\}$ mit $A \rightarrow B \rightarrow A$,
$\text{SCC}_2 = \{C,D\}$ mit $C \rightarrow D \rightarrow C$,
Querb\"{a}ngigkeit $A \rightarrow C$.
Korrekte Level: $\text{level}(C)=\text{level}(D)=0$, $\text{level}(A)=\text{level}(B)=1$.

Rekursiver Algorithmus, startend bei $A$:

\begin{verbatim}
level(A)  [Sentinel = 0]
  -> level(B) -> level(A): Zyklus, return 0  =>  B.level = 1
  -> level(C)  [Sentinel = 0]
       -> level(D) -> level(C): Zyklus, return 0  => D.level = 1
     C.level = 2   <- FALSCH (korrekt waere 0)
  A.level = max(1,2)+1 = 3   <- FALSCH (korrekt waere 1)
\end{verbatim}

\textbf{Satz:} Ein rekursiver Level-Algorithmus mit Sentinel-basierter
Zykluserkennung produziert f\"{u}r Graphen mit mehreren interagierenden SCCs
Ergebnisse, die von der Traversierungsreihenfolge abh\"{a}ngen.
Da die Traversierungsreihenfolge bei JVM-basierter Klassenladung der
Classpath-Reihenfolge folgt, sind die Ergebnisse \textbf{nicht deterministisch
\"{u}ber verschiedene Build-Konfigurationen}.

\subsection*{Die L\"{o}sung: zwei semantisch verschiedene Level-Konzepte}

Die Fehlermodi~1--3 haben eine gemeinsame Ursache: der Versuch,
\emph{architektonische Tiefe} und \emph{lokale Layoutposition} in einem einzigen
Level-Begriff zu vereinen. Die Pipeline trennt beide Konzepte konsequent:

\begin{itemize}
  \item \textbf{Phase~1} berechnet die \emph{architektonische Tiefe} --- den
        l\"{a}ngsten Pfad im globalen, zyklenfreien SCC-DAG \"{u}ber
        \textbf{alle Klassen des Projekts}.
  \item \textbf{Phase~2} berechnet die \emph{lokale Layout-Position} innerhalb
        eines Packages, beginnend bei~0, ausschlie{\ss}lich auf
        Geschwister-Abh\"{a}ngigkeiten.
\end{itemize}

\begin{verbatim}
Phase 1: Klassen-Graph -> Tarjan -> SCC-DAG (azyklisch)
           -> Longest-Path -> globale Level
           -> SCC-Equalisierung (alle Mitglieder = Maximum)
           -> Package-Level (Propagation nach oben)
           [Iterationsloop: Step 5 + Step 4b bis stabil]

Phase 2: Pro Package: Geschwister-Graph (nur interne Kanten)
           -> Tarjan -> lokaler SCC-DAG
           -> Longest-Path -> lokale Visualisierungs-Positionen
\end{verbatim}

\noindent\textbf{Hinweis:} Nach Phase~1 ist Level keine Eigenschaft einer einzelnen
Klasse, sondern des SCC-Knotens dem sie angeh\"{o}rt --- gemessen als l\"{a}ngster Pfad
im zyklenfreien Quotienten-Graphen (SCC-DAG) \"{u}ber alle Klassen des Projekts.

\begin{table}[h]
\small
\centering
\begin{tabular}{lll}
\hline
& \textbf{Phase 1} & \textbf{Phase 2} \\
\hline
Bedeutung    & Architekton. Tiefe im SCC-DAG & Layout-Position unter Geschwistern \\
Scope        & Global --- alle Klassen       & Lokal --- Kinder eines Packages \\
Startpunkt   & Blattklassen = 0              & Immer 0 f\"{u}r unterstes Element \\
Zyklen       & SCC-Kollaps vor Berechnung    & Lokaler Tarjan \\
Wertebereich & Hunderte m\"{o}glich           & Geschwisteranzahl $-$ 1 \\
\hline
\end{tabular}
\caption{Vergleich der beiden Level-Konzepte}
\end{table}

\subsection*{R4: eine verworfene Regel --- und was das lehrt}

Eine fr\"{u}he Fassung enthielt R4: \emph{,,Zyklisch voneinander abh\"{a}ngige
Kind-Packages m\"{u}ssen dasselbe Level teilen.``} Die Regel scheiterte in der Praxis:
Parent$\leftrightarrow$Child-Kanten, R\"{u}ckkanten des SCCBreakers und gemeinsame
Klassen-SCCs erzeugten Schein-Zyklen und f\"{a}lschliche Equalisierungen.
R4 wurde durch R2 ersetzt, das denselben Anspruch mit dem richtigen gefilterten
Graphen umsetzt.

\textbf{Die eigentliche Lehre:} Constraints f\"{u}r ein komplexes System korrekt zu
formulieren ist selbst schwer. Der Checker schafft einen \emph{doppelten Boden}:
er validiert nicht nur die Pipeline-Implementierung gegen die Regeln, sondern
deckt auch auf wenn die Regeln selbst zu grob sind. Dass R2 danach einen echten
Pipeline-Bug in \emph{software-ekg-7} fand, beweist dass die verschärfte Regel
pr\"{a}zise genug war um echte Fehler von Schein-Fehlern zu trennen.

\end{document}
